\documentclass[../DoAn.tex]{subfiles}
\begin{document}

\section{Đặt vấn đề}
\label{section:1.1}
Trong một trường đại học, phòng học và phòng họp là tài nguyên dùng chung có tần suất sử dụng cao. Ngoài thời khóa biểu cố định, giảng viên thường xuyên cần mượn thêm phòng để dạy bù, tổ chức seminar, họp bộ môn hoặc hướng dẫn nhóm đồ án. Với tòa nhà C7 mà nhóm khảo sát, mỗi tầng có bảy phòng bố trí theo sơ đồ chữ U, phòng số lẻ nằm một phía hành lang, phòng số chẵn nằm phía đối diện, xen giữa là thang máy, cầu thang và khu vệ sinh; phần lớn người dùng di chuyển bằng thang máy nên việc biết trước vị trí phòng trên từng tầng giúp tiết kiệm đáng kể thời gian đi lại.

Hiện nay quy trình mượn phòng chủ yếu dựa vào trao đổi trực tiếp, điện thoại hoặc tin nhắn với cán bộ cơ sở vật chất (CSVC). Cách làm này bộc lộ bốn hạn chế. Thứ nhất, người mượn khó biết phòng nào còn trống trong khung giờ mong muốn. Thứ hai, thông tin sức chứa và trang thiết bị như máy chiếu, micro, điều hòa không có sẵn để tra cứu nhanh. Thứ ba, người mượn không rõ phòng nằm ở phía nào của tầng. Thứ tư, các thông báo về kết quả duyệt, lịch nhận khóa hay thay đổi phòng không đến kịp thời. Ở phía quản lý, cán bộ CSVC phải tự đối chiếu lịch để tránh trùng, đồng thời xử lý hai phương thức tiếp cận phòng khác nhau là bàn giao thẻ từ và cấp mật khẩu (pass) cho khóa mã số.

Nếu bài toán được giải quyết, giảng viên có thể tự tra cứu và gửi yêu cầu trong vài thao tác, cán bộ CSVC có công cụ duyệt tập trung với quy tắc được kiểm tra tự động, và nhà trường có dữ liệu để đối soát lịch sử sử dụng phòng. Mô hình này cũng áp dụng được cho các tài nguyên đặt trước khác như phòng thí nghiệm, hội trường hoặc thiết bị dùng chung.

\section{Mục tiêu và phạm vi đề tài}
\label{section:1.2}
Các công cụ đang được dùng để quản lý lịch phòng có thể chia thành ba nhóm. Nhóm thứ nhất là sổ ghi chép hoặc bảng tính dùng chung, dễ triển khai nhưng không ngăn được trùng lịch và không phân quyền. Nhóm thứ hai là lịch trực tuyến như Google Calendar hay Microsoft Outlook với tài nguyên phòng; các công cụ này hiển thị được lịch bận nhưng không mô hình hóa quy tắc riêng của nhà trường như hạn đặt trước từ một đến ba ngày, giới hạn số yêu cầu đang hoạt động của mỗi người hoặc lịch bảo trì. Nhóm thứ ba là các phần mềm quản lý cơ sở vật chất thương mại, có đầy đủ chức năng nhưng chi phí cao và không gắn với loại khóa mà tòa nhà đang sử dụng. Điểm chung của các nhóm trên là tách rời việc \emph{đặt phòng} với việc \emph{tiếp cận phòng}: sau khi được duyệt, người mượn vẫn phải liên hệ riêng để nhận thẻ hoặc mật khẩu.

Xuất phát từ các hạn chế đó, bài tập lớn đặt mục tiêu xây dựng hệ thống RoomHub gồm ứng dụng di động cho người dùng cuối và dịch vụ nghiệp vụ Java. Phạm vi của đề tài bao gồm (i) đăng nhập, đăng ký và điều hướng theo hai vai trò Admin và User; (ii) tìm phòng trên sơ đồ tầng theo ngày giờ, sức chứa, thiết bị và loại khóa, rồi gửi yêu cầu ngay trên cùng màn hình; (iii) duyệt, từ chối và đổi phòng thay thế có kiểm tra điều kiện; (iv) quản lý tài khoản, phòng, lịch bảo trì và các tham số nghiệp vụ; (v) cấp quyền tiếp cận theo loại khóa, gồm hẹn lịch nhận khóa cơ hoặc thẻ từ và tạo mã PIN tạm thời cho khóa mã số; (vi) thông báo theo sự kiện. Các nghiệp vụ check-in bằng QR, xử lý no-show, audit log đầy đủ và báo cáo thống kê được mô tả trong quy trình nghiệp vụ nhưng nằm ngoài phạm vi phiên bản hiện tại.

\section{Định hướng giải pháp}
\label{section:1.3}
Nhóm tiếp cận bài toán theo phương pháp phân tích thiết kế hướng đối tượng. Phía máy chủ sử dụng Java 17 với Spring Boot 3.4 để xây dựng dịch vụ REST, Jakarta Persistence (JPA) để ánh xạ các lớp thực thể xuống cơ sở dữ liệu H2. Lựa chọn này phù hợp với yêu cầu của học phần Lập trình nâng cao về lập trình Java, đồng thời tận dụng cơ chế tiêm phụ thuộc, quản lý giao dịch và xử lý ngoại lệ tập trung của Spring. Phía người dùng là ứng dụng React Native viết bằng TypeScript để chạy trực tiếp trên điện thoại Android, nơi giảng viên thường thao tác khi di chuyển giữa các tầng.

Giải pháp tổ chức mã nguồn thành các module nghiệp vụ độc lập (xác thực, phòng, đặt phòng, kiểm soát truy cập, khóa thông minh, bảo trì, cấu hình, thông báo). Mọi quy tắc đặt phòng được đặt trong tầng dịch vụ và được kiểm tra lại ở máy chủ trong một giao dịch có khóa bi quan. Với phòng dùng khóa mã số, việc duyệt yêu cầu tự động cấp quyền cho đúng cặp người dùng -- phòng; người dùng tự tạo mã PIN sáu chữ số, ứng dụng gửi lệnh tới gateway để gateway mã hóa và phát bản tin MQTT tới khóa DLWA12 qua nền tảng OneIoT.

Kết quả đạt được là một sản phẩm chạy được trên thiết bị thật: APK Release độc lập không cần máy tính, mã nguồn di động khoảng 5.200 dòng TypeScript, dịch vụ Java 860 dòng với tám lớp thực thể, cùng bộ kiểm thử tự động cho các quy tắc chính. Các đóng góp nổi bật được trình bày chi tiết tại Chương~\ref{chapter:SolutionAndContribution}.

\section{Bố cục báo cáo}
\label{section:1.4}
Phần còn lại của báo cáo được tổ chức như sau.

Chương~\ref{chapter:Related_works} trình bày kết quả khảo sát hiện trạng mượn phòng và tài liệu quy trình nghiệp vụ, từ đó xây dựng biểu đồ use case tổng quát, các biểu đồ phân rã, quy trình nghiệp vụ đặt phòng, đặc tả năm use case quan trọng và các yêu cầu phi chức năng.

Chương~\ref{chapter:Methodology} giới thiệu nền tảng lý thuyết và công nghệ sử dụng, gồm lập trình hướng đối tượng với Java, Spring Boot và JPA, kiến trúc REST, React Native với TypeScript, giao thức MQTT và mã hóa AES; với mỗi công nghệ, báo cáo nêu yêu cầu mà công nghệ đó giải quyết và lý do lựa chọn so với phương án thay thế.

Chương~\ref{chapter:Experiment} trình bày thiết kế kiến trúc, thiết kế gói, thiết kế lớp, biểu đồ trình tự, thiết kế cơ sở dữ liệu và giao diện, sau đó mô tả kết quả xây dựng ứng dụng, các ca kiểm thử và mô hình triển khai trên thiết bị thật.

Chương~\ref{chapter:SolutionAndContribution} tập trung vào ba đóng góp chính: thuật toán kiểm tra xung đột lịch kết hợp khóa bi quan, mô hình cấp quyền tạo mã PIN theo phòng tích hợp khóa thông minh, và kiến trúc module hai chế độ dữ liệu.

Chương~\ref{chapter:conclusion} tổng kết kết quả, so sánh với các giải pháp hiện có, nêu hạn chế và hướng phát triển.

\end{document}
